The ablation study isolates the contributions of three core components of the graph-guided pipeline: (i) the section scheduler that determines generation order from the document graph, (ii) the enforcement of a strict JSON output schema for section-level generators, and (iii) the evidence-injection mechanism that supplies retrieved passages and explicit source pointers to the generator. For each component we compare the full graph-guided pipeline against a single-component ablation in which that component is disabled while all other components remain unchanged. The three ablation variants are implemented as follows.

\begin{itemize}
  \item \textbf{Scheduler ablation (No-Scheduler).} Generation proceeds in a linear, top-down order that ignores dependency edges in the document graph; retrieval and JSON constraints remain active. This variant isolates the scheduler's effect on citation allocation and coherence.
  \item \textbf{Schema ablation (Loose-JSON).} The generator is not required to produce outputs conforming to the strict JSON schema; instead, it may return unconstrained free-form text which is post-processed into the target fields. The scheduler and evidence injection are retained to measure the schema's effect on machine-parseability and citation-slot alignment.
  \item \textbf{Evidence ablation (No-Evidence).} Retrieved items and explicit source pointers are withheld from the generator; the scheduler and JSON constraints remain active. This variant isolates the contribution of direct evidence provisioning to citation coverage and factuality.
\end{itemize}

Evaluation for each variant uses the same test splits and annotation protocol described in Section~4.4. We measure (a) citation coverage (the fraction of substantive claims with at least one cited source), (b) hallucination rate (the fraction of claims judged inconsistent with available references), and (c) blind reviewer scores for overall quality. Differences between the full system and each ablation are reported as mean differences with standard deviations across the test set; statistical significance is assessed using paired tests and multiple-comparison correction as described below.

The following description of hypothesis-testing choices and thresholding is general background knowledge: we report two-sided paired tests (e.g., paired t-test or Wilcoxon signed-rank where normality is violated) and apply a Benjamini--Hochberg procedure to control the false discovery rate across the set of ablation comparisons.

We interpret ablation outcomes conservatively. An ablation that yields a reliably lower citation coverage or higher hallucination rate relative to the full pipeline is interpreted as evidence that the ablated component contributes to citation-aware factuality. Conversely, an ablation that primarily degrades blind reviewer scores but does not substantially change automated metrics suggests benefits that are stylistic or coherence-related rather than strictly evidential.

Limitations of this ablation analysis include sensitivity to retrieval quality and finite test-set size; these considerations affect the generality of inferences drawn from the comparisons and are discussed further in Section~6. Because run artifacts and numeric outputs are not included in this excerpt, the quantitative results and exact test statistics for each ablation are reported in Section~5.1 and in the supplementary materials.